\chapter{AI 深度学习全景族谱与零基础入门指南}

\begin{noteBox}[核心导读与背景]
\textbf{前言}：深度学习领域充斥着各种各样的“Net”（如 CNN、RNN、ResNet、Transformer、Diffusion），初学者极易迷失在浩瀚的术语海洋中。本篇为你梳理整张「AI 科技树」，讲清楚各种网络各司其职的关系，并给出最少弯路、最落地的零基础入门路线。
\end{noteBox}

\vspace{0.8em}\hrule\vspace{0.8em}

\section{[概览] 一、深度学习终极演化科技树（全景族谱）}

整个现代 AI 的发展，本质上只由两条主线驱动：\textbf{视觉线（看图像）} 和 \textbf{语言时序线（读文本）}。它们在经历各自数十年的探索后，最终被 \textbf{Transformer} 完成大一统！

\begin{figure}[htbp]
\centering
\begin{tikzpicture}[
  base/.style={draw=primary!80, fill=lightbg, rounded corners=2mm, align=center, font=\footnotesize\sffamily, inner sep=4pt, line width=0.8pt},
  vnode/.style={base, draw=secondary!80, fill=secondary!8, minimum width=3.8cm},
  nnode/.style={base, draw=accent!80, fill=accent!8, minimum width=3.8cm},
  tform/.style={base, draw=primary, fill=primary!15, line width=1.2pt, font=\bfseries\footnotesize\sffamily, minimum width=8cm},
  arr/.style={-{Stealth[scale=0.8]}, draw=gray!70, line width=0.8pt},
  darr/.style={-{Stealth[scale=0.8]}, dashed, draw=primary, line width=0.9pt}
]

% Level 0 & 1
\node[base, minimum width=6.5cm] (p) {1957 感知机 (Perceptron)\\单神经元，无法解决 XOR 异或};
\node[base, below=0.5cm of p, minimum width=6.5cm] (mlp) {1986 多层感知机 (MLP / BP)\\反向传播与现代神经网络基石};

% Level 2: CNN vs RNN
\node[vnode, below left=0.7cm and 0.4cm of mlp] (cnn) {1998 LeNet / CNN\\局部滑窗提取空间特征};
\node[nnode, below right=0.7cm and 0.4cm of mlp] (rnn) {1980s RNN (循环网络)\\时序备忘录处理文本};

% Level 3: AlexNet vs LSTM
\node[vnode, below=0.5cm of cnn] (alex) {2012 AlexNet (GPU 革命)\\算力与深度爆发};
\node[nnode, below=0.5cm of rnn] (lstm) {1997 LSTM / 2014 GRU\\门控机制缓解长程遗忘};

% Level 4: ResNet vs Seq2Seq
\node[vnode, below=0.5cm of alex] (resnet) {2015 ResNet (何恺明)\\残差连接 $x+F(x)$ 破百层};
\node[nnode, below=0.5cm of lstm] (seq2seq) {2014 Seq2Seq + Attention\\机器翻译学会“看重点”};

% Level 5: Transformer (Centered below both branches)
\node[tform, below=1.0cm of $(resnet.south)!0.5!(seq2seq.south)$] (trans) {2017 Transformer (Google 论文)\\彻底废黜循环时序！全矩阵并行点积！统一语言与多模态};

% Level 6: Successors
\node[base, below left=0.8cm and -0.8cm of trans, minimum width=3.2cm] (bert) {2018 BERT\\双向理解完形填空};
\node[base, below=0.8cm of trans, minimum width=3.6cm, fill=primary!10, draw=primary] (gpt) {2018--2023 GPT-1$\sim$4\\单向自回归接龙};
\node[base, below right=0.8cm and -0.8cm of trans, minimum width=3.2cm] (vit) {2020 ViT\\视觉 Transformer};

% Level 7: Modern LLM Era
\node[base, below=0.6cm of gpt, fill=accent!10, draw=accent, line width=1pt, font=\bfseries\footnotesize\sffamily, minimum width=8.5cm] (llm) {2023 至今：现代大模型时代 (Llama-3, Qwen-2.5, DeepSeek)\\核心转向：推理吞吐、显存架构优化、KV Cache、量化};

% Arrows
\draw[arr] (p) -- (mlp);
\draw[arr] (mlp) -| (cnn);
\draw[arr] (mlp) -| (rnn);
\draw[arr] (cnn) -- (alex);
\draw[arr] (alex) -- (resnet);
\draw[arr] (rnn) -- (lstm);
\draw[arr] (lstm) -- (seq2seq);

\draw[darr] (resnet.south) -- ++(0,-0.4) -| node[near start, left, font=\scriptsize] {残差连接思想} ($(trans.north west)+(0.8,0)$);
\draw[darr] (seq2seq.south) -- ++(0,-0.4) -| node[near start, right, font=\scriptsize] {注意力思想飞跃} ($(trans.north east)+(-0.8,0)$);

\draw[arr] (trans) -- (bert);
\draw[arr] (trans) -- (gpt);
\draw[arr] (trans) -- (vit);
\draw[arr] (gpt) -- (llm);

\end{tikzpicture}
\caption{AI 深度学习演化科技树全景族谱}
\end{figure}

\vspace{0.8em}\hrule\vspace{0.8em}

\section{[指南] 二、各种“Net”究竟都是干什么的？}

当你看到不同论文或文章里提到不同的 Net 时，可以对照下表快速定位：

\begin{center}
\small
\begin{tabularx}{\textwidth}{p{2.5cm} p{3.5cm} X X}
\toprule
\textbf{模型代号} & \textbf{全称} & \textbf{诞生初衷与解决的核心痛点} & \textbf{状态与历史地位} \\
\midrule
\textbf{MLP} & Multi-Layer Perceptron (多层感知机 / 全连接网络) & 最基础的万能函数拟合器，全连接矩阵乘法。但输入长度固定，无法直接处理大图像和长文章。 & 所有现代网络的基本积木（如 Transformer 里的 FFN） \\
\textbf{CNN} & Convolutional Neural Network (卷积神经网络) & \textbf{专为图像而生}。一张 1000x1000 的图片直接用全连接会导致参数爆炸，CNN 用一个小滑窗（卷积核）扫过图片提取边缘和纹理。 & 统治计算机视觉领域（2012-2020），经典代表：LeNet, AlexNet, VGG \\
\textbf{ResNet} & Residual Network (残差网络) & 网络堆得太深（超过 20 层）会导致梯度消失退化。何恺明发明\textbf{残差短路连接（$x + F(x)$）}，让网络轻松堆到 152 层甚至 1000 层！ & \textbf{现代所有大模型的标配机制}（Transformer 每层都有残差连接） \\
\textbf{RNN} & Recurrent Neural Network (循环神经网络) & \textbf{专为文本时序而生}。像穿糖葫芦一样逐字传递隐藏状态备忘录，解决文本长度不一的问题。 & 因串行无法利用 GPU 并行，且长程记忆衰减，已被淘汰 \\
\textbf{LSTM} & Long Short-Term Memory (长短期记忆网络) & 在 RNN 基础上引入遗忘门、输入门、输出门，缓解了 RNN 读到后面忘了前面的死穴。 & 在 2017 年前是机器翻译、语音识别的绝对霸主 \\
\textbf{Transformer} & 基于 Self-Attention 的全新网络架构 & \textbf{彻底废弃循环连接}，任意两词距离 $O(1)$，纯矩阵运算极致利用 GPU。不仅统领自然语言，还通过 ViT 杀入计算机视觉，通过 Sora/Diffusion 杀入视频生成。 & \textbf{当今全人类所有大模型的唯一通用底座} \\
\bottomrule
\end{tabularx}
\end{center}

\vspace{0.8em}\hrule\vspace{0.8em}

\section{[目标] 三、新手零基础想真正入门，该看什么？（绝不踩坑三部曲）}

网上的 AI 教程铺天盖地，90\% 都是死记硬背公式或调包调库，初学者往往陷入“懂了但又没完全懂”的痛苦中。

想要扎实、高效、不浪费时间地入门，\textbf{推荐且仅推荐以下三部封神级的黄金资源}：

\subsection{[第一阶段] 阶段 1：建立直观的数学与几何世界观（耗时约 2\textasciitilde{}3 天）}

\begin{itemize}
  \item \textbf{【绝对首选】3Blue1Brown 经典动画教程《深度学习的本质》（Bilibili / YouTube 免费看）}
  \item 搜索关键词：\texttt{3Blue1Brown 深度学习}
  \item 为什么神：全网最好的几何可视化，用生动的动态图形讲清楚：
  \item 什么是神经网络与激活函数？
  \item 什么是梯度下降（走下山坡）？
  \item 反向传播微积分链式法则的几何动画。
  \item 注意力机制（Attention）到底在几何空间中拉伸了什么。
  \item \textbf{建议}：把这 4 集视频当作看电影，看上两遍，你脑海里的数学直觉就彻底打通了！
\end{itemize}

\subsection{[第二阶段] 阶段 2：动手写代码与经典教科书（耗时约 1\textasciitilde{}2 周）}

\begin{itemize}
  \item \textbf{【教材推荐】李沐（前亚马逊首席科学家）团队《动手学深度学习》(Dive into Deep Learning / d2l.ai)}
  \item 中文官网完全免费在线阅读，配备完整可运行代码。
  \item 从基础张量、线性回归、MLP 讲到现代注意力机制，文字极度平实。
  \item \textbf{【硬核进阶】Andrej Karpathy（OpenAI 联合创始人、前特斯拉 AI 总监）的视频系列《Neural Networks: Zero to Hero》}
  \item 他在油管上从一个只有几行的微型框架 \texttt{micrograd} 开始，用纯 Python 逐行写出一个微型 GPT（nanoGPT）。
  \item 跟着他敲一遍，大模型的黑盒彻底透明。
\end{itemize}

\subsection{[第三阶段] 阶段 3：大模型推理与系统实战（面向生产与高性能）}

\begin{itemize}
  \item \textbf{【避坑指南：不要一上来就学“从零预训练大模型”】}
  \item 训练一个 70B 模型需要几千张英伟达 H100 显卡和几千万美元预算，对绝大多数工程师而言，学如何从头训练一个大模型完全脱离现实。
  \item \textbf{【正确的工程路径：深耕大模型推理（LLM Inference）】}
  \item 理解模型的前向传播逻辑与微观张量流（参见本项目的 \texttt{lab04\_transformer\_microscope}）；
  \item 理解为什么推理的核心瓶颈在于\textbf{显存容量与内存带宽（Memory Bandwidth）}（参见 \texttt{lab08\_memory\_and\_roofline}）；
  \item 深入业界工业标准框架（\textbf{vLLM}、\textbf{SGLang}、\textbf{llama.cpp}），学习它们的核心技术：\textbf{PagedAttention}、\textbf{连续批处理（Continuous Batching）}、\textbf{投机采样}、\textbf{量化}；
  \item 这也是当前 AI Infra、大模型私有化部署、高性能边缘计算最紧缺、含金量最高的技术方向！
\end{itemize}

\vspace{0.8em}\hrule\vspace{0.8em}

\begin{noteBox}[核心导读与背景]
{}[指南] \textbf{本仓库的完整路线}：上面三部曲是外部资源的推荐；本仓库自己的 14 篇理论（第 0\textasciitilde{}13 篇）+ 配套实验，以及它们与外部资源如何穿插学习，见 \textbf{LEARNING\_PATH.md}。
族谱里没有展开的细节——注意力在 2014 年的真正起源、从原版 Transformer 到 Llama/DeepSeek 的每一处改动——在 \textbf{第 5 篇}。
\end{noteBox}

\vspace{0.8em}\hrule\vspace{0.8em}

\textbf{下一篇}：\textbf{第 1 篇：计算机如何读懂语言？从人工规则到大模型演进史} ｜ \textbf{总路线}：\textbf{LEARNING\_PATH.md}
